\chapter*{Appendix G: Gauge Theory and Fiber--Bundle Structures in RSVP}
\addcontentsline{toc}{chapter}{Appendix G: Gauge Theory and Fiber--Bundle Structures in RSVP}

\section{Principal Bundles}

\subsection{Definition}

Let $G$ be a Lie group.  
A principal $G$-bundle consists of:
\[
\pi : P \to M,
\]
where $P$ is a smooth manifold equipped with a free, right $G$-action, such that $M = P/G$.

\begin{definition}[Local Trivialisation]
A neighbourhood $U\subseteq M$ admits a diffeomorphism
\[
\varphi : \pi^{-1}(U) \to U\times G
\]
that is $G$-equivariant.
\end{definition}

\section{Connections}

\subsection{Connection 1-Forms}

A connection on $P$ is a $\mathfrak{g}$-valued 1-form
\[
\omega \in \Omega^1(P,\mathfrak{g})
\]
satisfying:
\[
\omega(X^\#) = X, \qquad R_g^\ast \omega = \operatorname{Ad}_{g^{-1}} \omega,
\]
where $X^\#$ is the fundamental vector field associated to $X\in\mathfrak{g}$.

\subsection{Horizontal and Vertical Subspaces}

The tangent space splits:
\[
T_pP = \ker(\omega_p) \oplus \mathrm{Vert}_p(P).
\]

\section{Curvature}

\subsection{Curvature 2-Form}

The curvature of $\omega$ is
\[
\Omega = d\omega + \tfrac{1}{2}[\omega,\omega].
\]

\begin{definition}[Flat Connection]
A connection is flat if $\Omega = 0$.
\end{definition}

\section{Gauge Transformations}

\subsection{Gauge Group}

The gauge group is
\[
\mathcal{G} = \{ g : M\to G \ \text{smooth} \}.
\]

\subsection{Action on Connections}

\[
\omega \mapsto g^{-1}\omega g + g^{-1} dg.
\]

\section{Associated Bundles}

\subsection{Vector Bundles}

Given a representation $\rho : G\to \mathrm{GL}(V)$, the associated bundle is
\[
E = P \times_\rho V.
\]

Sections of $E$ correspond to $G$-equivariant functions $P\to V$.

\subsection{Connection on Associated Bundle}

A principal connection induces a covariant derivative on $E$:
\[
\nabla : \Gamma(E) \to \Gamma(E\otimes T^\ast M).
\]

\section{RSVP Field Bundles}

\subsection{Scalar Field}

The RSVP scalar field $\Phi$ is a section of the trivial real line bundle:
\[
\Phi \in \Gamma(M\times\mathbb{R}).
\]

\subsection{Vector Field}

The RSVP flow field $\mathbf{v}$ is a section of the tangent bundle:
\[
\mathbf{v} \in \Gamma(TM).
\]

\subsection{Entropy Field}

The entropy field $S$ is again scalar-valued:
\[
S \in \Gamma(M\times\mathbb{R}).
\]

\section{Gauge Freedom in RSVP}

\subsection{Curvature-Equivalent Representations}

Two field configurations $(\Phi,\mathbf{v},S)$ and $(\Phi',\mathbf{v}',S')$ are gauge-equivalent if related by:
\[
\Phi' = \Phi + f, \qquad \mathbf{v}' = \mathbf{v} + \nabla f, \qquad S' = S,
\]
for some scalar gauge potential $f$.

\subsection{Entropy Gauge}

Entropy admits gauge shifts preserving relative differences:
\[
S \mapsto S + c.
\]

\section{Bundle Interpretation of Polyxan Embeddings}

\subsection{Incidence Bundle}

To each hypergraph $\mathcal{G}$ there is a fiber bundle:
\[
\pi : \mathrm{Inc}(\mathcal{G}) \to |\mathcal{G}|.
\]

\subsection{Embedded Bundle}

An embedding $\iota : \mathrm{Inc}(\mathcal{G}) \to M$ is a bundle map respecting fibers.

\section{Curvature-Preserving Maps}

\begin{definition}[Gauge-Compatible Map]
A smooth map $F : M\to N$ is gauge-compatible when
\[
F^\ast \Omega_N = \Omega_M.
\]
\end{definition}

\begin{definition}[Curvature-Preserving Embedding]
An embedding $\iota$ is curvature-preserving if
\[
\iota^\ast \Omega = \Omega_{\mathcal{G}}
\]
for the curvature 2-form induced by the chosen connection on the incidence bundle.
\end{definition}

\section{Flatness and Semantic Stability}

Flatness of the induced connection corresponds to quine-stability:
\[
\Omega_{\mathcal{G}} = 0 \quad \Longleftrightarrow \quad \mathcal{G} \text{ is structurally stable under EHR rewriting}.
\]

\section{Gauge Fixing}

A gauge fixing is a global section of the principal bundle or a choice of representative in each gauge orbit.

Examples include:
\[
\nabla\cdot\mathbf{v} = 0, \qquad \Phi = \Phi_0, \qquad S = S_0.
\]

\section{Bundle Morphisms}

A morphism of RSVP bundles is a commutative diagram:
\[
\begin{array}{ccc}
P_1 & \xrightarrow{F} & P_2 \\
\downarrow &          & \downarrow \\
M_1 & \xrightarrow{f} & M_2
\end{array}
\]
with $F$ $G$-equivariant.

\section{Global Holonomy}

The holonomy group of $(P,\omega)$ is the subgroup of $G$ generated by parallel transport around closed loops in $M$.

\section{Holonomy and Semantic Curvature}

Holonomy characterises residual semantic curvature in embedding flows:
\[
\text{Flat holonomy} \quad \Longleftrightarrow \quad \text{zero semantic curvature}.
\]

\section{Reduction of Structure Group}

If the curvature constraints restrict $G$ to a subgroup $H$, there is a reduced bundle:
\[
P \to P/H.
\]

\section{Bundle Decomposition of RSVP Flows}

RSVP evolution can be expressed as:
\[
(\Phi_t,\mathbf{v}_t,S_t) = \nabla^\omega \sigma_t
\]
for a time-dependent section $\sigma_t$ of an associated RSVP bundle.

\section{Gauge-Invariant Quantities}

Gauge-invariant curvature scalars include:
\[
\|\Omega\|^2, \qquad \mathrm{Tr}(\Omega\wedge\Omega), \qquad |\nabla\Phi|^2.
\]

These correspond to RSVP invariants used in energy descent functionals.


